\chapter{Локальні аналітичні сховища (DWH)}

\section{Архітектура локального сховища}
Підсистема «ERP/1: Аналітика» призначена для виконання складних аналітичних запитів (OLAP) над даними транзакційних модулів (Mnesia, KVS) без перевантаження основного OLTP-контуру. Основна вимога до підсистеми --- робота on-premises, автономність, та повна інтеграція в екосистему Erlang/OTP без використання залежностей від Rust.

Архітектура сховища базується на використанні вбудованої стовпчикової бази даних DuckDB, що працює в ізольованому контурі. Дані з Mnesia трансформуються в оптимізовані стиснені файли Parquet та зберігаються у локальному S3-сумісному об'єктному сховищі (наприклад, Scality).

\section{DuckDB як вбудований аналітичний рушій}
DuckDB підключається безпосередньо до Erlang-вузла через C99 NIF (Native Implemented Functions) обгортку. Для запобігання блокування основних потоків планувальника Erlang (Erlang Schedulers), усі виклики DuckDB виконуються в контурі \textbf{Dirty CPU Schedulers}.

\subsection{Забезпечення ізоляції (Standalone)}
Ізоляція DuckDB від транзакційного ядра забезпечується шляхом:
\begin{itemize}
    \item Запуску DuckDB в окремому OS-процесі або ізольованій групі потоків (Thread Pool) з обмеженням CPU affinity;
    \item Обмеження виділеної оперативної пам'яті через конфігурацію DuckDB (\texttt{SET max\_memory = '8GB'});
    \item Використання ручного звільнення дескрипторів з'єднань за допомогою функцій очищення ресурсів Erlang NIF (\texttt{enif\_open\_resource\_type} та \texttt{duckdb\_close}).
\end{itemize}

\subsection{Горизонтальне масштабування (Scalability)}
Оскільки DuckDB є вбудованою однокористувацькою БД, горизонтальне масштабування аналітики реалізується через:
\begin{enumerate}
    \item \textbf{Федеративні SQL-запити (Federated Queries)} --- DuckDB здатна читати й об'єднувати гігабайти Parquet-файлів безпосередньо з віддаленого S3 за допомогою розширення \texttt{httpfs}. Це дозволяє виконувати розподілені обчислення над єдиним озером даних (Data Lake).
    \item \textbf{Peer-to-Peer реплікація} --- метадані та схеми даних синхронізуються між периферійними аналітичними вузлами через OTP-контур за допомогою легковагого протоколу консенсусу (Raft або Mnesia Sync).
\end{enumerate}

\section{Порівняльний аналіз аналітичних рішень}
У Додатку А викладено порівняльний аналіз трьох рішень для локальних інфраструктур DWH: ClickHouse, MonetDB та Apache Superset.

\begin{center}
\small
\begin{longtable}{p{3.5cm} p{3.5cm} p{3.5cm} p{3.5cm}}
\toprule
\textbf{Критерій} & \textbf{ClickHouse} & \textbf{MonetDB} & \textbf{Apache Superset} \\
\midrule
\textbf{Тип системи} & Розподілена стовпчикова СКБД & Вбудована / серверна C99 стовпчикова БД & Платформа BI та візуалізації звітів \\
\textbf{Мова реалізації} & C++17 / C++20 & Pure C (C99) & Python / React \\
\textbf{Рівень складності} & Високий (потребує окремого кластера) & Низький (легковажна стовпчикова СКБД) & Середній (потребує Python-середовища) \\
\textbf{Інтеграція з Erlang/OTP} & Через TCP протокол / HTTP API & Через C NIF або порт зв'язку & Окремий веб-сервер над аналітичними БД \\
\bottomrule
\end{longtable}
\end{center}

Завдяки реалізації MonetDB на чистому C99, вона є найбільш наближеним та легким закритим конкурентом ClickHouse для впровадження у периферійні вузли ERP/1, які мають жорсткі обмеження на ресурси.
